

\section{Цель работы}

Целью данной лабораторной работы является освоение методов принятия решений в условиях неопределенности. В рамках работы требуется применить классические критерии выбора оптимальной стратегии: критерий Вальда (максимин), критерий Сэвиджа (минимакс риска), критерий Гурвица, критерий Байеса-Лапласа и критерий Ходжа-Лемона.

Необходимо для каждого критерия вычислить оценочные значения $W_i$ по каждой стратегии, определить оптимальную стратегию и сравнить полученные результаты.

В соответствии с вариантом задания дана матрица выигрышей $A$ размерности $8 \times 10$ (8 стратегий, 10 состояний среды):

\[
A = \begin{pmatrix}
8 & 8 & 1 & 8 & 3 & 4 & 1 & 5 & 3 & 3 \\
3 & 6 & 6 & 2 & 5 & 6 & 3 & 3 & 6 & 6 \\
8 & 5 & 6 & 4 & 2 & 1 & 5 & 4 & 3 & 2 \\
7 & 3 & 8 & 8 & 8 & 8 & 4 & 1 & 6 & 8 \\
3 & 6 & 7 & 3 & 1 & 6 & 8 & 8 & 6 & 8 \\
9 & 9 & 7 & 9 & 1 & 1 & 6 & 6 & 4 & 6 \\
3 & 5 & 8 & 3 & 3 & 5 & 2 & 6 & 1 & 7 \\
8 & 4 & 4 & 8 & 5 & 7 & 4 & 5 & 1 & 8
\end{pmatrix}
\]

\clearpage







\section{Теоретическая часть}

Задача принятия решения в условиях неопределенности возникает в тех случаях, когда лицо, принимающее решение (ЛПР), не располагает информацией о вероятностях возможных состояний внешней среды. В отличие от задач в условиях риска, где вероятности состояний известны или могут быть оценены, здесь приходится выбирать стратегию, опираясь лишь на матрицу выигрышей и собственные представления о разумном поведении. Формально задача формулируется следующим образом.

Имеется $m$ возможных стратегий (альтернатив) $W_1, W_2, \ldots, W_m$ и $n$ возможных состояний среды (природы) $S_1, S_2, \ldots, S_n$, вероятности которых неизвестны. Задана матрица выигрышей $A = (a_{ij})$, где $a_{ij}$ -- выигрыш (полезность, доход, эффект) при выборе стратегии $W_i$ и состоянии среды $S_j$. Требуется выбрать стратегию, которая является оптимальной в некотором разумном смысле.

Поскольку понятие оптимальности в условиях неопределенности не является однозначным, в теории принятия решений разработан ряд критериев, каждый из которых отражает определённую позицию ЛПР по отношению к неопределенности и риску. Одни критерии ориентированы на наихудший исход (пессимистический подход), другие -- на наилучший (оптимистический), третьи пытаются найти компромисс между ними или опираются на предположение о равновероятности состояний среды. Рассмотрим основные из них.

\textbf{Критерий Вальда} (максиминный критерий, критерий крайнего пессимизма) ориентирован на наихудшее состояние среды. ЛПР предполагает, что природа будет действовать против него, и старается обеспечить себе максимальный гарантированный результат. Для каждой стратегии определяется её минимальный выигрыш:
\[
W_i = \min_{j} a_{ij}.
\]
Оптимальной считается стратегия, обеспечивающая максимальное значение этого минимума:
\[
W = \max_{i} \min_{j} a_{ij}.
\]
Критерий Вальда гарантирует, что выигрыш при любом состоянии среды будет не меньше найденного значения $W$. Однако он чрезмерно пессимистичен и может приводить к выбору заведомо невыгодных стратегий, если среди состояний среды есть маловероятные, но крайне неблагоприятные.

\textbf{Критерий Сэвиджа} (минимаксный критерий риска) также является пессимистическим, но учитывает не абсолютные выигрыши, а величину упущенной выгоды. Используется матрица рисков $R = (r_{ij})$, элементы которой вычисляются по формуле:
\[
r_{ij} = \max_{k} a_{kj} - a_{ij}.
\]
Величина $r_{ij}$ показывает, насколько выигрыш при выборе стратегии $W_i$ в состоянии $S_j$ меньше максимально возможного выигрыша в этом состоянии. Для каждой стратегии определяется максимальный риск:
\[
W_i = \max_{j} r_{ij}.
\]
Оптимальной считается стратегия, минимизирующая этот максимум:
\[
W = \min_{i} \max_{j} r_{ij}.
\]
Критерий Сэвиджа позволяет избежать больших потерь, но, как и критерий Вальда, ориентирован на наихудший случай.

\textbf{Критерий Гурвица} (критерий пессимизма-оптимизма) является взвешенной комбинацией наилучшего и наихудшего исходов и позволяет учесть склонность ЛПР к риску. При параметре $\lambda \in [0,1]$:
\[
W_i = \lambda \min_{j} a_{ij} + (1-\lambda) \max_{j} a_{ij}.
\]
Параметр $\lambda$ называется коэффициентом пессимизма. При $\lambda = 1$ критерий превращается в критерий Вальда, при $\lambda = 0$ -- в максимаксный критерий (крайнего оптимизма). При $\lambda = 1/2$ критерий превращается в усреднение минимального и максимального выигрышей. Оптимальной считается стратегия с максимальным значением $W_i$:
\[
W = \max_{i} W_i.
\]
Критерий Гурвица позволяет ЛПР регулировать степень своей осторожности, выбирая соответствующее значение $\lambda$.

\textbf{Критерий Байеса-Лапласа} предполагает, что все состояния среды равновероятны, то есть $p_j = 1/n$. Это предположение носит характер «принципа недостаточного основания» Лапласа: если нет информации о вероятностях, все состояния считаются одинаково возможными. Для каждой стратегии вычисляется средний ожидаемый выигрыш:
\[
W_i = \frac{1}{n} \sum_{j=1}^{n} a_{ij}.
\]
Оптимальной считается стратегия с максимальным значением:
\[
W = \max_{i} W_i.
\]
Критерий Байеса-Лапласа является наиболее оптимистичным из рассмотренных, поскольку ориентируется на средний результат и не учитывает экстремальные исходы.

\textbf{Критерий Ходжа-Лемона} (обобщённый критерий Гурвица) при параметре $\nu \in [0,1]$ учитывает как средний ожидаемый выигрыш, так и минимальный выигрыш, то есть объединяет подходы критериев Байеса-Лапласа и Вальда:
\[
W_i = \nu \cdot \frac{1}{n} \sum_{j=1}^{n} a_{ij} + (1-\nu) \cdot \min_{j} a_{ij}.
\]
Параметр $\nu$ отражает степень доверия ЛПР к предположению о равновероятности состояний среды. При $\nu = 1$ критерий совпадает с критерием Байеса-Лапласа, при $\nu = 0$ -- с критерием Вальда. При $\nu = 1/2$ критерий усредняет среднее арифметическое строки и минимальный элемент строки. Оптимальной считается стратегия с максимальным значением $W_i$.

Таким образом, каждый из рассмотренных критериев даёт свою рекомендацию, и окончательный выбор стратегии остаётся за ЛПР, которое может учитывать как формальные расчёты, так и дополнительные содержательные соображения. Сравнение результатов, полученных по разным критериям, позволяет получить более полную картину задачи и оценить устойчивость выбора.









\clearpage

\section{Практическая часть}

Критерий Вальда:

Для каждой строки найдём минимальный элемент:
\begin{align*}
W_1 &= \min(8,8,1,8,3,4,1,5,3,3) = 1, \\
W_2 &= \min(3,6,6,2,5,6,3,3,6,6) = 2, \\
W_3 &= \min(8,5,6,4,2,1,5,4,3,2) = 1, \\
W_4 &= \min(7,3,8,8,8,8,4,1,6,8) = 1, \\
W_5 &= \min(3,6,7,3,1,6,8,8,6,8) = 1, \\
W_6 &= \min(9,9,7,9,1,1,6,6,4,6) = 1, \\
W_7 &= \min(3,5,8,3,3,5,2,6,1,7) = 1, \\
W_8 &= \min(8,4,4,8,5,7,4,5,1,8) = 1.
\end{align*}

Максимальное значение среди полученных:
\[
W = \max(1,2,1,1,1,1,1,1) = 2 \quad \text{при стратегии } W_2.
\]

Критерий Сэвиджа:

Сначала найдём максимумы по каждому столбцу:
\[
\begin{gathered}
\max_j a_{1j} = 9, \quad \max_j a_{2j} = 9, \quad \max_j a_{3j} = 8, \quad \max_j a_{4j} = 9, \quad \max_j a_{5j} = 8, \\
\max_j a_{6j} = 8, \quad \max_j a_{7j} = 8, \quad \max_j a_{8j} = 8, \quad \max_j a_{9j} = 6, \quad \max_j a_{10j} = 8.
\end{gathered}
\]

Построим матрицу рисков $R = (r_{ij})$, где $r_{ij} = \max_k a_{kj} - a_{ij}$:

\[
R = \begin{pmatrix}
1 & 1 & 7 & 1 & 5 & 4 & 7 & 3 & 3 & 5 \\
6 & 3 & 2 & 7 & 3 & 2 & 5 & 5 & 0 & 2 \\
1 & 4 & 2 & 5 & 6 & 7 & 3 & 4 & 3 & 6 \\
2 & 6 & 0 & 1 & 0 & 0 & 4 & 7 & 0 & 0 \\
6 & 3 & 1 & 6 & 7 & 2 & 0 & 0 & 0 & 0 \\
0 & 0 & 1 & 0 & 7 & 7 & 2 & 2 & 2 & 2 \\
6 & 4 & 0 & 6 & 5 & 3 & 6 & 2 & 5 & 1 \\
1 & 5 & 4 & 1 & 3 & 1 & 4 & 3 & 5 & 0
\end{pmatrix}
\]

Для каждой строки найдём максимальный риск:
\begin{align*}
W_1 &= \max(1,1,7,1,5,4,7,3,3,5) = 7, \\
W_2 &= \max(6,3,2,7,3,2,5,5,0,2) = 7, \\
W_3 &= \max(1,4,2,5,6,7,3,4,3,6) = 7, \\
W_4 &= \max(2,6,0,1,0,0,4,7,0,0) = 7, \\
W_5 &= \max(6,3,1,6,7,2,0,0,0,0) = 7, \\
W_6 &= \max(0,0,1,0,7,7,2,2,2,2) = 7, \\
W_7 &= \max(6,4,0,6,5,3,6,2,5,1) = 6, \\
W_8 &= \max(1,5,4,1,3,1,4,3,5,0) = 5.
\end{align*}

Минимальное значение среди полученных:
\[
W = \min(7,7,7,7,7,7,6,5) = 5 \quad \text{при стратегии } W_8.
\]

Критерий Гурвица (при $\lambda = 1/2$):

Для каждой строки вычислим:
\[
W_i = \frac{1}{2} \min_j a_{ij} + \frac{1}{2} \max_j a_{ij}.
\]

\begin{align*}
W_1 &= \frac{1 + 8}{2} = 4.5, \\
W_2 &= \frac{2 + 6}{2} = 4, \\
W_3 &= \frac{1 + 8}{2} = 4.5, \\
W_4 &= \frac{1 + 8}{2} = 4.5, \\
W_5 &= \frac{1 + 8}{2} = 4.5, \\
W_6 &= \frac{1 + 9}{2} = 5, \\
W_7 &= \frac{1 + 8}{2} = 4.5, \\
W_8 &= \frac{1 + 8}{2} = 4.5.
\end{align*}

Максимальное значение:
\[
W = \max(4.5, 4, 4.5, 4.5, 4.5, 5, 4.5, 4.5) = 5 \quad \text{при стратегии } W_6.
\]

Критерий Байеса--Лапласа:

В предположении равновероятности всех состояний ($p_j = 1/10$) вычислим средний выигрыш для каждой стратегии:
\[
W_i = \frac{1}{10} \sum_{j=1}^{10} a_{ij}.
\]

\begin{align*}
W_1 &= \frac{8+8+1+8+3+4+1+5+3+3}{10} = \frac{44}{10} = 4.4, \\
W_2 &= \frac{3+6+6+2+5+6+3+3+6+6}{10} = \frac{46}{10} = 4.6, \\
W_3 &= \frac{8+5+6+4+2+1+5+4+3+2}{10} = \frac{40}{10} = 4, \\
W_4 &= \frac{7+3+8+8+8+8+4+1+6+8}{10} = \frac{61}{10} = 6.1, \\
W_5 &= \frac{3+6+7+3+1+6+8+8+6+8}{10} = \frac{56}{10} = 5.6, \\
W_6 &= \frac{9+9+7+9+1+1+6+6+4+6}{10} = \frac{58}{10} = 5.8, \\
W_7 &= \frac{3+5+8+3+3+5+2+6+1+7}{10} = \frac{43}{10} = 4.3, \\
W_8 &= \frac{8+4+4+8+5+7+4+5+1+8}{10} = \frac{54}{10} = 5.4.
\end{align*}

Максимальное значение:
\[
W = \max(4.4, 4.6, 4, 6.1, 5.6, 5.8, 4.3, 5.4) = 6.1 \quad \text{при стратегии } W_4.
\]

Критерий Ходжа--Лемона (при $\nu = 1/2$):

Для каждой строки вычислим:
\[
W_i = \frac{1}{2} \cdot \frac{1}{n} \sum_{j=1}^{n} a_{ij} + \frac{1}{2} \cdot \min_j a_{ij}.
\]

Используя ранее найденные средние и минимумы:
\begin{align*}
W_1 &= \frac{4.4 + 1}{2} = 2.7, \\
W_2 &= \frac{4.6 + 2}{2} = 3.3, \\
W_3 &= \frac{4.0 + 1}{2} = 2.5, \\
W_4 &= \frac{6.1 + 1}{2} = 3.55, \\
W_5 &= \frac{5.6 + 1}{2} = 3.3, \\
W_6 &= \frac{5.8 + 1}{2} = 3.4, \\
W_7 &= \frac{4.3 + 1}{2} = 2.65, \\
W_8 &= \frac{5.4 + 1}{2} = 3.2.
\end{align*}

Максимальное значение:
\[
W = \max(2.7, 3.3, 2.5, 3.55, 3.3, 3.4, 2.65, 3.2) = 3.55 \quad \text{при стратегии } W_4.
\]


Решения игры следующие:
\begin{itemize}
    \item по критерию Вальда: $2$ при стратегии $2$;
    \item по критерию Сэвиджа: $5$ при стратегии $8$;
    \item по критерию Гурвица: $5$ при стратегии $6$;
    \item по критерию Байеса--Лапласа: $6.1$ при стратегии $4$;
    \item по критерию Ходжа--Лемона: $3.55$ при стратегии $4$.
\end{itemize}

\clearpage

\section{Заключение}

В ходе выполнения данной лабораторной работы были освоены методы принятия решений в условиях неопределенности. Для заданной матрицы выигрышей размерности $8 \times 10$ были применены пять классических критериев: Вальда, Сэвиджа, Гурвица, Байеса--Лапласа и Ходжа--Лемона. Для каждого критерия вычислены оценочные значения $W_i$ по всем восьми стратегиям и определена оптимальная стратегия.

Анализ результатов показывает, что разные критерии дают различные рекомендации, что является ожидаемым, поскольку каждый критерий отражает определённую позицию лица, принимающего решения, по отношению к неопределенности. Критерий Вальда, ориентированный на крайний пессимизм, выбирает стратегию $W_2$ с гарантированным выигрышем $2$. Критерий Сэвиджа, минимизирующий максимальный риск, указывает на стратегию $W_8$. Критерий Гурвица при $\lambda=1/2$ рекомендует стратегию $W_6$. Критерии Байеса--Лапласа и Ходжа--Лемона, учитывающие средние ожидаемые значения, оба указывают на стратегию $W_4$, причём стратегия $W_4$ имеет наибольший средний выигрыш ($6.1$) и наибольшее значение по критерию Ходжа--Лемона ($3.55$).

Таким образом, стратегия $W_4$ представляется наиболее предпочтительной с точки зрения среднего ожидаемого выигрыша, в то время как стратегия $W_2$ обеспечивает наилучший гарантированный результат. Выбор окончательной стратегии зависит от склонности лица, принимающего решения, к риску и от конкретных условий задачи. Цель лабораторной работы была полностью достигнута.

\clearpage

\section{Список литературы}

\begin{enumerate}
    \item Саати Т. Принятие решений при зависимостях и обратных связях: Аналитические сети. -- М.: Книжный дом «ЛИБРОКОМ», 2009. -- 360 с.
    \item Теория принятия решений в 2 т. Том 2: учебник и практикум для вузов / В. Г. Халин [и др.]; ответственный редактор В. Г. Халин. -- Москва: Издательство Юрайт, 2020. -- 431 с.
    \item Баранов, В. В. Процессы принятия управляющих решений, мотивированных интересами: учебное пособие / В. В. Баранов. -- Москва: ФИЗМАТЛИТ, 2005. -- 296 с.
    \item Петросян Л.А., Зенкевич Н.А., Шевкопляс Е.В. Теория игр. -- СПб.: БХВ-Петербург, 2012. -- 432 с.
    \item Лубенец Ю.В. Теория принятия решений: Учеб. пособие. -- Липецк: Изд-во ЛГТУ, 2011. -- 80 с.
\end{enumerate}
